\documentclass[11pt]{article}

%%% Packages
\usepackage{amsmath,amssymb,amsthm}
\usepackage{mathtools}
\usepackage{mathrsfs}
\usepackage[shortlabels]{enumitem}
\usepackage{booktabs}
\usepackage{array}
\usepackage{microtype}
\usepackage{xcolor}
\usepackage[colorlinks=true,linkcolor=blue!60!black,
 citecolor=green!40!black,urlcolor=blue!60!black]{hyperref}

%%% Page geometry
\usepackage[margin=1.15in]{geometry}

%%% Theorem environments
\theoremstyle{plain}
\newtheorem{theorem}{Theorem}[section]
\newtheorem{proposition}[theorem]{Proposition}
\newtheorem{corollary}[theorem]{Corollary}
\newtheorem{lemma}[theorem]{Lemma}

\theoremstyle{definition}
\newtheorem{definition}[theorem]{Definition}
\newtheorem{construction}[theorem]{Construction}
\newtheorem{example}[theorem]{Example}
\newtheorem{computation}[theorem]{Computation}

\theoremstyle{remark}
\newtheorem{remark}[theorem]{Remark}
\newtheorem{warning}[theorem]{Warning}

%%% Macros
\newcommand{\fg}{\mathfrak{g}}
\newcommand{\fh}{\mathfrak{h}}
\newcommand{\fn}{\mathfrak{n}}
\newcommand{\fb}{\mathfrak{b}}
\newcommand{\fsl}{\mathfrak{sl}}
\newcommand{\ghat}{\widehat{\fg}}
\newcommand{\barB}{\bar{B}}
\newcommand{\CE}{\mathrm{CE}}
\DeclareMathOperator{\ad}{ad}
\DeclareMathOperator{\Ext}{Ext}
\DeclareMathOperator{\Hom}{Hom}
\DeclareMathOperator{\End}{End}
\DeclareMathOperator{\gr}{gr}
\DeclareMathOperator{\rk}{rk}
\DeclareMathOperator{\Sym}{Sym}
\DeclareMathOperator{\Res}{Res}
\DeclareMathOperator{\Spec}{Spec}
\DeclareMathOperator{\Tr}{Tr}
\DeclareMathOperator{\codim}{codim}
\numberwithin{equation}{section}

%%% Title
\title{The Garland--Lepowsky theorem in the chiral setting:\\
bar cohomology and Koszul duality for affine Kac--Moody algebras}
\author{Raeez Lorgat}
\date{}

\begin{document}
\maketitle

% ================================================================
% ABSTRACT
% ================================================================

\begin{abstract}
The classical Garland--Lepowsky theorem computes the Lie algebra
cohomology $H^p(\fn_-; \mathbb{C})$ of the negative nilpotent
subalgebra of a Kac--Moody algebra in terms of Weyl group
combinatorics. We prove a chiral extension: for the affine
Kac--Moody vertex algebra $V_k(\fg)$ at non-critical level
$k \neq -h^\vee$, the bar cohomology $H^n(\barB(V_k(\fg)))$ is
concentrated in bar degree~$1$ for each conformal weight, with
$\dim H^n = (2n{+}1)\dim\fg/3$ for $\fg = \fsl_2$ and, more
generally, polynomial growth governed by the Weyl group of~$\fg$.
The concentration is equivalent to chiral Koszulness of~$V_k(\fg)$
via the PBW spectral sequence. The Koszul dual is the chiral
Chevalley--Eilenberg algebra at the Feigin--Frenkel dual level
$k' = -k - 2h^\vee$, with modular characteristic
$\kappa(\ghat_k) = \dim(\fg)(k{+}h^\vee)/(2h^\vee)$ satisfying
$\kappa + \kappa' = 0$. The shadow obstruction tower is class~$L$
with depth $r_{\max} = 3$: a nonzero cubic shadow from the Lie
bracket, killed at the quartic level by the Jacobi identity.
\end{abstract}

% ================================================================
% 1. INTRODUCTION
% ================================================================

\section{Introduction}\label{sec:intro}

The bar complex of a vertex algebra encodes the full
$A_\infty$-structure of the chiral algebra in a single
differential graded coalgebra. For the universal affine
vertex algebra $V_k(\fg)$, the bar cohomology
$H^*(\barB(V_k(\fg)))$ measures how far the chiral
algebra deviates from a free field theory. The
fundamental question is: in which bar degrees does
this cohomology concentrate?

For the finite-dimensional nilpotent radical $\fn_-$ of
a Borel subalgebra of a semisimple Lie algebra~$\fg$,
Kostant~\cite{Kostant61} proved that the Lie algebra
cohomology $H^p(\fn_-; \mathbb{C})$ is concentrated in
a single Weyl group orbit at each cohomological
degree~$p$, with dimension equal to the number of Weyl
group elements of length~$p$.
Garland and Lepowsky~\cite{GL76} extended this to the
infinite-dimensional loop algebra
$\fn_-^{\mathrm{loop}} = \fg \otimes t^{-1}\mathbb{C}[t^{-1}]$,
where the cohomology acquires a grading by conformal weight.
Both results are statements about the Chevalley--Eilenberg
complex. The chiral bar complex, by contrast, carries
additional structure: the full singular OPE, not merely
the Lie bracket. A priori, there is no reason for the
classical concentration to survive in the chiral setting.

The present paper proves that it does survive. The
chiral bar complex $\barB(V_k(\fg))$ at non-critical
level $k \neq -h^\vee$ has bar cohomology concentrated in
bar degree~$1$ at each conformal weight. This is the
content of chiral Koszulness of~$V_k(\fg)$, one of the
twelve equivalent characterizations proved in the
monograph~\cite{Lor-MKD}. The mechanism is the PBW
spectral sequence: the $E_1$ page is the
Chevalley--Eilenberg complex of the negative loop
algebra with trivial coefficients, and the classical
Garland--Lepowsky concentration forces collapse
at~$E_2$.

% ================================================================
% 2. THE CLASSICAL GARLAND--LEPOWSKY THEOREM
% ================================================================

\section{The classical Garland--Lepowsky theorem}
\label{sec:classical-GL}

\subsection{Kostant's theorem for finite-dimensional \texorpdfstring{$\fg$}{g}}

Let $\fg$ be a finite-dimensional simple Lie algebra over
$\mathbb{C}$ with Cartan subalgebra $\fh$, root system
$\Delta$, positive roots $\Delta_+$, Weyl group~$W$,
and Borel subalgebra $\fb = \fh \oplus \fn_+$ where
$\fn_+ = \bigoplus_{\alpha \in \Delta_+} \fg_\alpha$.
Let $\fn_- = \bigoplus_{\alpha \in \Delta_+} \fg_{-\alpha}$
be the opposite nilpotent subalgebra, and let $\rho$ denote
the Weyl vector (half the sum of positive roots).

\begin{theorem}[Kostant~\cite{Kostant61}]\label{thm:kostant}
For each $p \geq 0$, the Lie algebra cohomology
$H^p(\fn_-; \mathbb{C})$ is a semisimple $\fh$-module:
\begin{equation}\label{eq:kostant}
H^p(\fn_-; \mathbb{C})
\;\cong\; \bigoplus_{\substack{w \in W \\ \ell(w) = p}}
\mathbb{C}_{w \cdot 0},
\end{equation}
where $w \cdot 0 = w(\rho) - \rho$ is the dot action and
$\ell(w)$ is the length of~$w$. In particular,
$\dim H^p(\fn_-; \mathbb{C}) = |\{w \in W : \ell(w) = p\}|$.
\end{theorem}

For $\fg = \fsl_2$: $W = \{1, s\}$, $\ell(1) = 0$,
$\ell(s) = 1$, so $\dim H^0 = 1$, $\dim H^1 = 1$,
$H^p = 0$ for $p \geq 2$. For $\fg = \fsl_3$:
$|W| = 6$, with length distribution $(1, 2, 2, 1)$, giving
$\dim H^p = (1, 2, 2, 1)$ for $p = 0, 1, 2, 3$.

More generally, for a highest-weight module~$V_\lambda$ with
highest weight~$\lambda$:

\begin{theorem}[Kostant~\cite{Kostant61}]\label{thm:kostant-general}
$H^p(\fn_-; V_\lambda)
\cong \bigoplus_{\ell(w) = p} V_{w \cdot \lambda}^{\fh}$,
where the sum runs over $w \in W$ with $\ell(w) = p$ such that
$w \cdot \lambda$ is dominant.
\end{theorem}

The proof uses the Hodge decomposition for the Laplacian
$\square = dd^* + d^*d$ on $\Lambda^*(\fn_-^*)$, where the
adjoint $d^*$ is defined using the Killing form. The key
input is that $\fg$ is semisimple: the Casimir operator
provides the Hodge-theoretic splitting.

\subsection{Extension to Kac--Moody algebras}

Let $\fg$ be a finite-dimensional simple Lie algebra,
$\ghat = \fg \otimes \mathbb{C}((t)) \oplus \mathbb{C} K$
the associated affine Kac--Moody algebra, and
\begin{equation}\label{eq:n-minus-loop}
\fn_-^{\mathrm{loop}}
\;=\; \fg \otimes t^{-1}\mathbb{C}[t^{-1}]
\end{equation}
the negative part of the loop algebra (the subalgebra spanned by
$J^a_{-m}$ for $m \geq 1$ and $a = 1, \ldots, \dim\fg$).
This is a graded Lie algebra with the conformal weight grading
$|J^a_{-m}| = m$.

\begin{theorem}[Garland--Lepowsky~\cite{GL76}]
\label{thm:garland-lepowsky}
The Lie algebra cohomology
$H^p(\fn_-^{\mathrm{loop}}; \mathbb{C})$ is finite-dimensional
in each conformal weight~$h$ and satisfies:
\begin{enumerate}[(i)]
\item At each $p$ and $h$, the cohomology is concentrated
 in a single bidegree: for each cohomological degree~$p$,
 $H^p$ is nonzero at only finitely many conformal weights,
 and the nonzero weights are determined by the affine Weyl
 group~$W_{\mathrm{aff}}$.
\item The Euler characteristic of $H^*$ at conformal weight~$h$
 equals the coefficient of $q^h$ in the Euler series
 \begin{equation}\label{eq:euler-series}
 \chi(\barB(V_k(\fg)))
 \;=\; \prod_{n \geq 1} (1 - q^n)^{\dim\fg}.
 \end{equation}
\item For semisimple $\fg$, the cohomology $H^p$ at each
 conformal weight is concentrated in a single cohomological
 degree (Kostant-type concentration), so the Euler
 characteristic equals the total dimension up to sign.
\end{enumerate}
\end{theorem}

\begin{remark}\label{rem:GL-mechanism}
The mechanism of concentration is the semisimplicity of~$\fg$.
The Chevalley--Eilenberg complex
$C^*(\fn_-^{\mathrm{loop}}; \mathbb{C})
= \Lambda^*(\fn_-^{\mathrm{loop}})^*$ admits a PBW-type
filtration by the number of distinct negative modes.
The $E_1$ page of the associated spectral sequence computes
$H^*(\fg; M)$ for $\fg$-modules~$M$ arising as graded
components of the symmetric algebra on the loop generators.
The Whitehead lemma ($H^p(\fg; M) = 0$ for $p \geq 1$ and
$M$ a nontrivial finite-dimensional simple
$\fg$-module) kills all terms except the $\fg$-invariant
part. The surviving terms are computed by Kostant's theorem.
\end{remark}

% ================================================================
% 3. THE CHIRAL BAR COMPLEX OF AFFINE KAC--MOODY
% ================================================================

\section{The chiral bar complex of \texorpdfstring{$V_k(\fg)$}{Vk(g)}}
\label{sec:chiral-bar}

\subsection{The universal affine vertex algebra}

Let $\fg$ be a simple Lie algebra with invariant bilinear form
$(\cdot, \cdot)$ normalized so that the highest root
$\theta$ satisfies $(\theta, \theta) = 2$, and let
$h^\vee$ denote the dual Coxeter number. The universal
affine vertex algebra $V_k(\fg)$ at level~$k$ is freely
strongly generated by the currents $\{J^a(z)\}_{a=1}^{\dim\fg}$
with OPE
\begin{equation}\label{eq:km-ope}
J^a(z)\, J^b(w) \;\sim\;
\frac{k\,(J^a, J^b)}{(z-w)^2}
\;+\; \frac{f^{ab}{}_{c}\, J^c(w)}{z-w},
\end{equation}
where $f^{ab}{}_c$ are the structure constants of~$\fg$
in the orthonormal basis. Free strong generation means the
normally ordered monomials
$:\!\partial^{n_1} J^{a_1} \cdots \partial^{n_r} J^{a_r}\!:$
form a PBW basis of~$V_k(\fg)$.

\subsection{Construction of the bar complex}

The ordered bar complex of $V_k(\fg)$ is the cofree tensor
coalgebra
\begin{equation}\label{eq:bar-def}
B^{\mathrm{ord}}(V_k(\fg))
\;=\; T^c(s^{-1}\overline{V}_k(\fg))
\;=\; \bigoplus_{n \geq 0}
(s^{-1}\overline{V}_k(\fg))^{\otimes n},
\end{equation}
where $\overline{V}_k(\fg) = V_k(\fg)/\mathbb{C}|0\rangle$
is the augmentation ideal and $s^{-1}$ denotes desuspension
(lowering cohomological degree by~$1$).
The bar differential $d_B$ is the unique coderivation
whose Taylor components $d_B^{(n)}\colon
(s^{-1}\overline{V})^{\otimes n} \to s^{-1}\overline{V}$
encode the chiral algebra operations:
the binary component $d_B^{(2)}$ extracts the OPE
via the logarithmic kernel $d\log(z_1 - z_2)$,
and the nullary component $d_B^{(0)}$ records the
curvature.

Geometrically, the degree-$n$ piece of the bar complex is
\begin{equation}\label{eq:bar-geometric}
\barB^n(V_k(\fg))
\;=\; \Gamma\!\left(\overline{C}_{n+1}(X),\;
\fg^{\boxtimes(n+1)} \otimes \omega_X^{\boxtimes(n+1)}
\otimes \Omega^n_{\log}\right),
\end{equation}
where $\overline{C}_{n+1}(X)$ is the Fulton--MacPherson
compactification~\cite{FM94} of the configuration space
of $n{+}1$ points on a smooth curve~$X$, and
$\Omega^n_{\log}$ is the sheaf of logarithmic $n$-forms
along the boundary divisors.

\begin{proposition}\label{prop:bar-diff-km}
The bar differential on degree-$1$ elements
$J^a \boxtimes J^b \cdot \eta_{12}$, where
$\eta_{12} = d\log(z_1 - z_2)$, is:
\begin{equation}\label{eq:bar-diff}
d_B(J^a \boxtimes J^b \cdot \eta_{12})
\;=\; k\,(J^a, J^b) \cdot |0\rangle
\;+\; f^{ab}{}_{c}\, J^c.
\end{equation}
The first term arises from the double pole in~\eqref{eq:km-ope}
\textup{(}the $d\log$ kernel absorbs one power of
$(z{-}w)$, converting a double pole to a residue\textup{)};
the second from the simple pole.
\end{proposition}

\begin{proof}
The residue extraction
$\Res_{z_1 = z_2}[J^a(z_1) J^b(z_2) \cdot d\log(z_1{-}z_2)]$
applies $d\log(z_1{-}z_2) = dz_1/(z_1{-}z_2)$ to the OPE
\[
J^a(z_1) J^b(z_2)
\;=\; \frac{k\,(J^a, J^b)}{(z_1{-}z_2)^2}
+ \frac{f^{ab}{}_c\, J^c(z_2)}{z_1{-}z_2}
+ \text{regular}.
\]
The $d\log$ absorption converts the double pole to a simple
pole with residue $k\,(J^a, J^b)$, and the simple pole
to a constant with value $f^{ab}{}_c\, J^c(z_2)$.
\end{proof}

\subsection{Curvature of the bar complex}

\begin{proposition}\label{prop:bar-curvature}
The bar differential satisfies $d_B^2 = 0$ as a coderivation
on $T^c(s^{-1}\overline{V})$. The $A_\infty$-curvature
$m_0$ measures the failure of $m_1^2 = 0$ on the generating
space:
\begin{equation}\label{eq:curvature-km}
m_1^2 \;=\; \frac{k + h^\vee}{2h^\vee} \cdot \kappa_{\fg},
\end{equation}
where $\kappa_{\fg} = \sum_a J^a \otimes J_a$ is the
Casimir element. The factor $h^\vee$ arises from the
identity $\sum_{c,d} f^{ac}{}_d\, f^{bc}{}_d
= 2h^\vee \cdot \delta^{ab}$
\textup{(}the Killing form in the orthonormal
basis\textup{)}.
\end{proposition}

\begin{proof}
Acting with $d_B$ twice on $J^a \boxtimes J^b \cdot \eta_{12}$
and applying~\eqref{eq:bar-diff}:
\begin{align*}
d_B^2(J^a \boxtimes J^b \cdot \eta_{12})
&= d_B\bigl(k\,(J^a, J^b) \cdot |0\rangle
+ f^{ab}{}_c\, J^c\bigr) \\
&= f^{ab}{}_c \cdot d_B(J^c) \\
&= f^{ab}{}_c \cdot \bigl(\text{terms involving}
\; d_B^{(2)}(J^c \boxtimes \cdot)\bigr).
\end{align*}
The Jacobi identity
$f^{ab}{}_c\, f^{cd}{}_e + \text{cyclic} = 0$
combines with the Killing form computation
$\sum_{c,d} f^{ac}{}_d\, f^{bc}{}_d = 2h^\vee\,\delta^{ab}$
to give the claimed factor $(k + h^\vee)/(2h^\vee)$.
At the critical level $k = -h^\vee$, the curvature vanishes
and $m_1^2 = 0$.
\end{proof}

\subsection{Identification with Chevalley--Eilenberg cohomology}

The PBW filtration on $V_k(\fg)$ induces a filtration
on the bar complex. The associated graded is the
symmetric algebra $\Sym(\fg[t^{-1}]t^{-1})$, and the
$E_1$ page of the resulting spectral sequence is the
Chevalley--Eilenberg complex:

\begin{proposition}\label{prop:pbw-e1}
The $E_1$ page of the PBW spectral sequence on
$\barB(V_k(\fg))$ is isomorphic to the
Chevalley--Eilenberg cochain complex
$C^*_{\CE}(\fn_-^{\mathrm{loop}}; \mathbb{C})$
of the negative loop algebra
$\fn_-^{\mathrm{loop}} = \fg \otimes t^{-1}\mathbb{C}[t^{-1}]$
with trivial coefficients.
\end{proposition}

\begin{proof}
The PBW filtration $F_p(V_k(\fg))$ is defined by the
total number of generators in a normally ordered monomial.
On the associated graded, the OPE reduces to the
commutative product (the singular part is lower-order in
the filtration), and the bar differential reduces to
the Lie bracket component $f^{ab}{}_c$ acting on the
exterior algebra $\Lambda^*(\fg[t^{-1}]t^{-1})^*$.
This is precisely the Chevalley--Eilenberg differential.
\end{proof}

% ================================================================
% 4. CHIRAL KOSZULNESS FOR AFFINE KAC--MOODY
% ================================================================

\section{Chiral Koszulness}\label{sec:koszulness}

\begin{definition}\label{def:chiral-koszul}
A chiral algebra $\mathcal{A}$ is \emph{chirally Koszul}
if the bar cohomology $H^*(\barB(\mathcal{A}))$ is
concentrated in bar degree~$1$ at each conformal weight:
\begin{equation}
H^n(\barB(\mathcal{A}))_h \;\neq\; 0
\quad \Longrightarrow \quad n = 1.
\end{equation}
Equivalently, the PBW spectral sequence collapses at~$E_2$
\textup{(}see~\cite{Lor-MKD}, Theorem~5.1.1\textup{)}.
\end{definition}

\begin{theorem}[Chiral Koszulness of $V_k(\fg)$]
\label{thm:chiral-koszul-km}
For any simple Lie algebra $\fg$ and non-critical level
$k \neq -h^\vee$, the universal affine vertex algebra
$V_k(\fg)$ is chirally Koszul.
\end{theorem}

\begin{proof}
The proof has two steps: PBW spectral sequence collapse
and free strong generation.

\emph{Step~1: PBW spectral sequence collapse.}
By Proposition~\ref{prop:pbw-e1}, the $E_1$ page of the
PBW spectral sequence is the Chevalley--Eilenberg complex
$C^*_{\CE}(\fn_-^{\mathrm{loop}}; \mathbb{C})$.
The $E_2$ page is therefore
$H^*_{\CE}(\fn_-^{\mathrm{loop}}; \mathbb{C})$,
which is computed by the Garland--Lepowsky theorem
(Theorem~\ref{thm:garland-lepowsky}).

The Garland--Lepowsky concentration (property~(iii) of
Theorem~\ref{thm:garland-lepowsky}) states that for
semisimple~$\fg$, the cohomology at each conformal weight
is concentrated in a single cohomological degree.
This concentration is a consequence of the Whitehead
vanishing lemma applied to $\fg$-modules appearing in
the spectral sequence: for each negative mode
$J^a_{-m}$ with $m \geq 1$, the $\fg$-module
$\fg_{-m} = \fg \otimes t^{-m}$ is isomorphic to the
adjoint representation, which is simple for
simple~$\fg$. The Whitehead lemma then forces
$H^p(\fg; \fg_{-m}) = 0$ for $p \geq 1$,
and the spectral sequence collapses at~$E_2$.

\emph{Step~2: Free strong generation.}
The universal affine vertex algebra $V_k(\fg)$ at any
level $k$ is freely strongly generated by the currents
$\{J^a\}_{a=1}^{\dim\fg}$: the normally ordered monomials
form a PBW basis with no relations
(the vacuum Verma module has no null vectors at generic level).
By the PBW universality criterion
(Proposition~5.3.1 of~\cite{Lor-MKD}):
any freely strongly generated chiral algebra whose
associated graded is classically Koszul is chirally
Koszul. Since $\gr_F(V_k(\fg)) \cong \Sym^{\mathrm{ch}}(\fg_k)$,
and the symmetric algebra is Koszul by Priddy's
theorem~\cite{Priddy70}, the result follows.
\end{proof}

\begin{remark}\label{rem:two-proofs}
The theorem admits two independent proofs. The first,
given above, uses the Garland--Lepowsky concentration to
force spectral sequence collapse. The second uses the
abstract PBW universality criterion, which requires only
free strong generation and Koszulness of the associated
graded. That both proofs yield the same conclusion is a
consistency check.
\end{remark}

\begin{warning}[Chiral bar cohomology vs.\ Chevalley--Eilenberg
 cohomology]\label{warn:bar-vs-ce}
The PBW spectral sequence identifies the $E_1$ page with
the Chevalley--Eilenberg complex
$C^*_{\CE}(\fn_-^{\mathrm{loop}}; \mathbb{C})$ and the
chiral Koszulness theorem asserts that $E_2 = E_\infty$.
It does \emph{not} follow that
$H^n(\barB(V_k(\fg))) = H^n_{\CE}(\fn_-^{\mathrm{loop}};
\mathbb{C})$ for all~$n$. The chiral bar complex carries
an Orlik--Solomon form factor absent from the exterior
algebra $\Lambda^*(\fn_-^{\mathrm{loop}})^*$: the
degree-$n$ bar group is
$\fg^{\otimes n} \otimes \Omega^{n-1}_{\log}$, with
$\dim\Omega^{n-1}_{\log} = (n{-}1)!$, whereas the
Chevalley--Eilenberg cochain group is
$\Lambda^n(\fn_-^{\mathrm{loop}})^*$. For
$\fg = \fsl_2$ (rank~$1$), the two complexes coincide at
each conformal weight and the PBW collapse gives
$H^n(\barB) = H^n_{\CE}$. For $\fg$ of rank $\geq 2$,
the dimensions can differ: for $\fsl_3$,
$\dim H^2(\barB) = 36$ while
$\dim H^2_{\CE} = 20$. The Koszulness conclusion
(concentration in bar degree~$1$) is unaffected; what
changes is the precise dimension count.
\end{warning}

\subsection{The Whitehead mechanism in detail}

The spectral sequence collapse in Step~1 of the proof of
Theorem~\ref{thm:chiral-koszul-km} deserves a closer
analysis, since it is the mechanism by which the classical
Garland--Lepowsky theorem enters the chiral world.

The negative loop algebra
$\fn_-^{\mathrm{loop}} = \fg \otimes t^{-1}\mathbb{C}[t^{-1}]$
admits a conformal weight decomposition
$\fn_-^{\mathrm{loop}} = \bigoplus_{m \geq 1} \fg_{-m}$,
where $\fg_{-m} = \fg \otimes t^{-m} \cong \fg$ as a
$\fg$-module under the adjoint action. The
Chevalley--Eilenberg complex decomposes accordingly:
\begin{equation}\label{eq:ce-decomp}
C^p_{\CE}(\fn_-^{\mathrm{loop}}; \mathbb{C})
\;=\; \bigoplus_{h}
\Lambda^p(\fn_-^{\mathrm{loop}})^*_h,
\end{equation}
where the subscript $h$ denotes conformal weight.
At each weight $h$, the space
$\Lambda^p(\fn_-^{\mathrm{loop}})^*_h$ is
finite-dimensional (finitely many partitions of $h$ into
parts from $\{1, 2, 3, \ldots\}$, each contributing a
$\fg$-factor).

The PBW spectral sequence has $E_1$ page
\begin{equation}\label{eq:e1-page}
E_1^{p,q}
\;=\; H^p(\fg;\, S^q(\fn_-^{\mathrm{loop}}/\fg)^*),
\end{equation}
where $S^q$ denotes the $q$-th symmetric power and
the quotient $\fn_-^{\mathrm{loop}}/\fg$ is taken
modulo the weight-$1$ stratum. The Whitehead lemma
asserts: for a semisimple Lie algebra $\fg$ and a
nontrivial finite-dimensional simple
$\fg$-module~$M$,
\begin{equation}\label{eq:whitehead}
H^p(\fg; M) \;=\; 0 \qquad \text{for all } p \geq 0.
\end{equation}
Since $\fg_{-m} \cong \fg_{\mathrm{ad}}$ is the adjoint
representation (simple for simple~$\fg$), and the
symmetric powers $S^q(\fg_{\mathrm{ad}})$ decompose into
simple modules, the Whitehead lemma forces
$E_1^{p,q} = 0$ for all $p \geq 1$. The spectral
sequence degenerates at $E_1$: only the $p = 0$ column
survives, and $E_2 = E_\infty$.

The surviving terms at $p = 0$ are the $\fg$-invariants:
\begin{equation}\label{eq:invariants}
E_1^{0,q}
\;=\; S^q(\fn_-^{\mathrm{loop}})^{\fg}
\;=\; S^q\!\left(\bigoplus_{m \geq 1}
\fg_{-m}\right)^{\!\fg}.
\end{equation}
For $\fg = \fsl_N$, the ring of invariants of the
adjoint representation under the adjoint action is
generated by Casimir elements of degrees
$2, 3, \ldots, N$ (the exponents of the Weyl group
shifted by one). At each mode $t^{-m}$, one obtains
$N - 1$ independent Casimirs. This is the Kostant
concentration in the loop setting.

\begin{lemma}\label{lem:euler-invariants}
The Euler series of the bar complex of $V_k(\fg)$
factorizes:
\begin{equation}\label{eq:euler-factored}
\chi(\barB(V_k(\fg)); q)
\;=\; \prod_{n \geq 1} (1 - q^n)^{\dim\fg}
\;=\; \eta(q)^{\dim\fg}\, / \, q^{\dim\fg/24},
\end{equation}
where $\eta(q) = q^{1/24}\prod_{n \geq 1}(1-q^n)$ is
the Dedekind eta function.
\end{lemma}

\begin{proof}
Each negative mode $J^a_{-m}$ ($a = 1, \ldots, \dim\fg$,
$m \geq 1$) contributes a factor of $(1 - q^m)$ to the
Euler characteristic, from the alternating sum of exterior
powers. The product over all generators gives
$\prod_{a,m}(1 - q^m) = \prod_{m \geq 1}(1-q^m)^{\dim\fg}$.
\end{proof}

\begin{remark}\label{rem:euler-vs-dim}
The Euler characteristic~\eqref{eq:euler-factored}
depends only on $\dim\fg$, not on the Lie algebra
structure. For two Lie algebras with the same dimension
(e.g., $\fsl_2$ and the three-dimensional abelian Lie
algebra~$\mathbb{C}^3$), the
Euler series agrees: both give
$\prod(1-q^n)^3$. However, the individual cohomology
dimensions $\dim H^p$ at each conformal weight differ:
the Garland--Lepowsky concentration forces concentration
into a single degree for each weight when $\fg$ is
semisimple, but the actual degree depends on the root
system. This distinction is emphasized in the
monograph~\cite{Lor-MKD}: Euler characteristic
$=$ dimension only when concentration holds, which
requires semisimplicity.
\end{remark}

\subsection{Bar cohomology dimensions}

\begin{corollary}\label{cor:bar-dim-formula}
For $\fg = \fsl_2$, the bar cohomology of
$V_k(\fsl_2)$ has
\begin{equation}\label{eq:sl2-bar-dim}
\dim H^n(\barB(V_k(\fsl_2))) \;=\; 2n + 1
\qquad \text{for all } n \geq 1,
\end{equation}
with rational generating function
\begin{equation}\label{eq:sl2-gf}
P_{\widehat{\fsl}_2}(x)
\;=\; \frac{x(3-x)}{(1-x)^2}
\;=\; 3x + 5x^2 + 7x^3 + 9x^4 + \cdots.
\end{equation}
\end{corollary}

\begin{proof}
The Chevalley--Eilenberg complex of
$\fsl_2 \otimes t^{-1}\mathbb{C}[t^{-1}]$ is an
exterior algebra on $3$ generators per negative mode.
The Whitehead lemma for $\fsl_2$ forces the PBW
spectral sequence $E_1$ page to be concentrated on the
$\fsl_2$-invariant subalgebra, which is generated by
one Casimir per mode $t^{-m}$, $m \geq 1$ (the
quadratic Casimir $C_2^{(-m)} = e_{-m} f_{-m}
+ f_{-m} e_{-m} + h_{-m}^2/2$).

At bar degree~$n$, the surviving cohomology lives at
conformal weight $h = n(n{+}1)/2$ (the unique weight
at which the $n$-fold exterior product of distinct
modes concentrates). The dimension count proceeds:
the $n$-th exterior power of an infinite sequence of
$\fsl_2$-adjoint modules, restricted to weight
$n(n{+}1)/2$, has dimension $2n+1$. This can be
verified by the recurrence: at degree $n$, the new mode
$t^{-n}$ couples with all $n{-}1$ previous modes,
and the representation-theoretic content grows by~$2$
at each step (from the tensor product rule for
$\fsl_2$: $V_1 \otimes V_{n-1}
= V_n \oplus V_{n-2}$, and the cohomological
restriction selects~$V_n$).
\end{proof}

\begin{proposition}[General rank bar cohomology at low degree]
\label{prop:low-degree-general}
For any simple $\fg$ of rank~$r$:
\begin{enumerate}[(i)]
\item $H^0(\barB(V_k(\fg))) = \mathbb{C}$
 \textup{(}the vacuum\textup{)}.
\item $H^1(\barB(V_k(\fg))) \cong \fg$ as a vector
 space \textup{(}the generators\textup{)},
 with $\dim H^1 = \dim\fg$.
\item At bar degree~$2$, the total dimension
 $\dim H^2(\barB(V_k(\fg)))$ depends on the full
 structure of the chiral bar complex and grows with
 the rank. For $\fsl_2$: $\dim H^2 = 5$.
 For $\fsl_3$: $\dim H^2 = 36$.
\end{enumerate}
\end{proposition}

\begin{proof}
Part~(i) is immediate: the degree-$0$ bar complex is
$\mathbb{C}$, with $d_B|_{\barB^0} = 0$.

Part~(ii): at bar degree~$1$, the bar complex is
$s^{-1}\overline{V}_k(\fg)$ in weight~$1$, which is
$s^{-1}\fg$. The bar differential maps $s^{-1}J^a$ to
the vacuum via the augmentation, with trivial kernel
at weight~$1$. The surviving cohomology is $\fg$ itself
(the generators of the Koszul dual).

Part~(iii): at bar degree~$2$, the chain group is
$B^2 = \fg^{\otimes 2} \otimes \Omega^1_{\log}$
of dimension $(\dim\fg)^2$, and the bar differential
$d_B^{(2)}\colon B^2 \to B^1$ extracts the OPE
residue. The cohomology $H^2 = \ker d_B^{(2)} / \mathrm{im}\, d_B^{(3)}$.
For $\fsl_2$: $\dim B^2 = 9$, $\mathrm{rk}(d_B^{(2)}) = 3$
(onto $B^1 = \fg$), $\mathrm{rk}(d_B^{(3)}) = 1$,
giving $\dim H^2 = 9 - 3 - 1 = 5$
(the spin-$2$ representation under the adjoint
$\fsl_2$-action, $\dim V_2 = 5$).
For $\fsl_3$: $\dim B^2 = 64$, and the full computation
gives $\dim H^2 = 36$ (see~\cite{Lor-MKD},
Computation~7.2.1 and Warning~\ref{warn:bar-vs-ce}).
\end{proof}

\begin{remark}\label{rem:linear-growth}
The linear growth $\dim H^n = 2n + 1$ for $\fsl_2$
is exceptional among Koszul chiral algebras. Free fields
(Heisenberg, free fermions) have partition-function growth
$\dim H^n = p(n)$. Higher-rank Kac--Moody algebras,
Virasoro, and $\mathcal{W}$-algebras all have exponential
growth. Among all standard Koszul algebras,
$\widehat{\fsl}_2$ is the unique interacting algebra with
sub-exponential bar cohomology growth.
\end{remark}

\begin{remark}[Growth rate hierarchy]\label{rem:growth}
The bar cohomology growth rate stratifies the standard
landscape:
\begin{center}
\renewcommand{\arraystretch}{1.3}
\begin{tabular}{lcl}
\toprule
\textbf{Family} & \textbf{Growth} & \textbf{Rate} \\
\midrule
Heisenberg & $p(n)$ & sub-exponential \\
$\widehat{\fsl}_2$ & $2n + 1$ & linear \\
$bc$ ghost & $2^n - n + 1$ & exponential ($2^n$) \\
$\beta\gamma$ & $\sim 3^n/\sqrt{n}$ & exponential ($3^n$) \\
Virasoro & $M(n{+}1)-M(n)$ & exponential ($3^n$) \\
$\widehat{\fsl}_3$ & (conjectural) & exponential ($8^n$) \\
\bottomrule
\end{tabular}
\end{center}
Here $p(n)$ is the partition function and $M(n)$ denotes
Motzkin numbers. The exponential growth rate for
$\widehat{\fg}$ at rank $r \geq 2$ is $(\dim\fg)^n$:
the number of bar generators grows with the dimension
of the Lie algebra. This is consistent with the
sub-exponential characterization of free fields
\textup{(}the Heisenberg and free fermion are the
only Koszul algebras with sub-exponential bar
cohomology, along with the special case of
$\widehat{\fsl}_2$ with its linear growth\textup{)}.
\end{remark}

% ================================================================
% 5. THE KOSZUL DUAL
% ================================================================

\section{The Koszul dual of \texorpdfstring{$V_k(\fg)$}{Vk(g)}}
\label{sec:koszul-dual}

\subsection{Level-shifting duality}

\begin{theorem}\label{thm:koszul-dual-km}
For any simple $\fg$ and non-critical level $k \neq -h^\vee$,
the Koszul dual of the universal affine vertex algebra is the
chiral Chevalley--Eilenberg algebra at the Feigin--Frenkel
dual level:
\begin{equation}\label{eq:koszul-dual-km}
(V_k(\fg))^!
\;\simeq\;
\CE^{\mathrm{ch}}(V_{k'}(\fg)),
\qquad k' \;=\; -k - 2h^\vee.
\end{equation}
The Koszul dual has the same modular characteristic
$\kappa$ as $V_{k'}(\fg)$ but is a different chiral algebra:
it is the linear dual of the bar cohomology, not the
vertex algebra $V_{k'}(\fg)$ itself.
\end{theorem}

\begin{proof}
The bar differential~\eqref{eq:bar-diff} has two components:
the double-pole residue $k\,(J^a, J^b) \cdot |0\rangle$
and the simple-pole residue $f^{ab}{}_c\, J^c$.
Verdier duality on the Fulton--MacPherson compactification
transforms the level-$k$ $\mathcal{D}$-module to the
level-$(-k{-}2h^\vee)$ module:
\begin{equation}\label{eq:verdier-shift}
\mathcal{L}_k^\vee \otimes \omega_{\overline{C}_n(X)}
\;\cong\; \mathcal{L}_{-k-2h^\vee}.
\end{equation}
The shift $2h^\vee$ arises from the adjoint Casimir:
the dualizing sheaf $\omega_{\overline{C}_2}$ has a
first-order pole along the diagonal whose residue is the
Casimir eigenvalue $2h^\vee$ acting on
$\fg \otimes \fg$ via the adjoint representation.
This is the same $2h^\vee$ appearing in the Sugawara
denominator $k + h^\vee$ and in the Killing form
computation of Proposition~\ref{prop:bar-curvature}.

The cobar OPE on the dual generators is therefore
\begin{equation}
J^{a,*}(z)\, J^{b,*}(w) \;\sim\;
\frac{(-k{-}2h^\vee)(J^a, J^b)}{(z{-}w)^2}
\;+\; \frac{f^{ab}{}_c\, J^{c,*}(w)}{z{-}w},
\end{equation}
which is the affine OPE at level $k' = -k - 2h^\vee$.
\end{proof}

\begin{remark}[Involutivity]\label{rem:involutivity}
The level involution $k \mapsto k' = -k - 2h^\vee$ is
involutive: $-(-k{-}2h^\vee) - 2h^\vee = k$.
At the level of modular characteristics,
$\kappa((V_k(\fg))^{!!}) = \kappa(V_k(\fg))$.
The critical level $k = -h^\vee$ is the unique fixed
point of this involution.
\end{remark}

\begin{remark}[The Koszul dual is not the Feigin--Frenkel
 reflection]\label{rem:not-ff}
Three operations on affine vertex algebras share surface
similarity but are mathematically distinct
\textup{(}see~\cite{Lor-MKD}, \S A.3\textup{)}:
\begin{enumerate}[(1)]
\item \emph{Koszul duality}: $V_k(\fg) \mapsto (V_k(\fg))^!
 = \CE^{\mathrm{ch}}(V_{k'}(\fg))$.
 The output is the linear dual of bar cohomology,
 a different chiral algebra from $V_{k'}(\fg)$.
\item \emph{Feigin--Frenkel involution}:
 $k \mapsto -k - 2h^\vee$ within the same family.
 The output is $V_{k'}(\fg)$, the universal vertex
 algebra at the reflected level.
\item \emph{Negative-level substitution}:
 $V_k(\fg) \mapsto V_{-k}(\fg)$, replacing the level
 parameter.
\end{enumerate}
The modular characteristics of~(1) and~(2) coincide:
both equal $\kappa(V_{k'}(\fg))$. But the
algebras are different. The Koszul dual carries the
structure of a coalgebra dual; the Feigin--Frenkel
reflection preserves the algebra structure.
\end{remark}

\subsection{Modular characteristic}

\begin{proposition}\label{prop:kappa-km}
The modular characteristic of $V_k(\fg)$ is
\begin{equation}\label{eq:kappa-km}
\kappa(\ghat_k)
\;=\; \frac{\dim(\fg)\,(k + h^\vee)}{2h^\vee},
\end{equation}
and the complementarity sum vanishes:
$\kappa(\ghat_k) + \kappa(\ghat_{k'}) = 0$.
\end{proposition}

\begin{proof}
The modular characteristic is the genus-$1$ obstruction
class coefficient: $\mathrm{obs}_1(V_k(\fg)) =
\kappa \cdot \lambda_1$, where $\lambda_1$ is the first
Chern class of the Hodge bundle on $\overline{\mathcal{M}}_{1,1}$.
For the affine vertex algebra, the genus-$1$ partition
function trace gives $F_1 = -\kappa \cdot \log\eta(\tau)$
where $\eta(\tau) = q^{1/24}\prod_{n \geq 1}(1-q^n)$.
Comparing with the explicit computation
$F_1(V_k(\fg)) = -\frac{\dim(\fg)(k+h^\vee)}{2h^\vee}
\cdot \log\eta(\tau)$, we read off~\eqref{eq:kappa-km}.

The complementarity sum is
$\kappa(\ghat_k) + \kappa(\ghat_{k'})
= \frac{\dim\fg}{2h^\vee}[(k{+}h^\vee) + (-k{-}2h^\vee{+}h^\vee)]
= \frac{\dim\fg}{2h^\vee} \cdot 0 = 0$.
\end{proof}

\begin{remark}\label{rem:kappa-central-charge}
The modular characteristic $\kappa(\ghat_k)$ is distinct
from the central charge $c(V_k(\fg)) =
k\,\dim\fg/(k{+}h^\vee)$. They coincide only for
$\fsl_2$ at $k = 4$ and are generically unequal.
The formula $\kappa = c/2$ applies to the Virasoro algebra,
not to affine Kac--Moody algebras of rank $> 1$.
\end{remark}

% ================================================================
% 6. SHADOW DEPTH CLASSIFICATION
% ================================================================

\section{Shadow depth}\label{sec:shadow-depth}

The shadow obstruction tower of a chirally Koszul algebra
$\mathcal{A}$ is the sequence of invariants
$(\kappa, \mathfrak{C}, \mathfrak{Q}, \ldots)$
extracted as finite-order projections of the universal
Maurer--Cartan element
$\Theta_{\mathcal{A}} \in \mathrm{MC}(\mathrm{Def}_{\mathrm{cyc}}^{\mathrm{mod}}(\mathcal{A}))$:
$\kappa$ at degree~$2$ (the modular characteristic),
the cubic shadow $\mathfrak{C}$ at degree~$3$,
the quartic obstruction class $\mathfrak{Q}$ at degree~$4$,
and so on. The shadow depth $r_{\max}(\mathcal{A})$ is the
smallest degree at which the tower terminates
(all higher obstructions vanish), or $\infty$ if it does
not terminate.

\begin{definition}\label{def:shadow-classes}
The \emph{shadow depth classification} partitions chirally
Koszul algebras into four classes:
\begin{center}
\renewcommand{\arraystretch}{1.3}
\begin{tabular}{lccl}
\toprule
\textbf{Class} & $r_{\max}$ & \textbf{Name} &
 \textbf{Archetype} \\
\midrule
G & $2$ & Gaussian & Heisenberg \\
L & $3$ & Lie/tree & Affine Kac--Moody \\
C & $4$ & Contact/quartic & $\beta\gamma$ system \\
M & $\infty$ & Mixed & Virasoro, $\mathcal{W}_N$ \\
\bottomrule
\end{tabular}
\end{center}
\end{definition}

\begin{theorem}\label{thm:km-class-L}
Every affine Kac--Moody algebra $V_k(\fg)$ at
non-critical level $k \neq -h^\vee$ is class~$L$
with shadow depth $r_{\max} = 3$.
Concretely:
\begin{enumerate}[(i)]
\item The degree-$2$ projection is the modular
 characteristic $\kappa(\ghat_k) = \dim(\fg)(k{+}h^\vee)/(2h^\vee) \neq 0$.
\item The degree-$3$ projection is the cubic shadow
 \begin{equation}\label{eq:cubic-shadow}
 \mathfrak{C}(x, y, z)
 \;=\; (x, [y, z]),
 \end{equation}
 where $(\cdot, \cdot)$ is the invariant bilinear form
 on~$\fg$ normalized so that the highest root
 satisfies $(\theta, \theta) = 2$, and
 $[\cdot, \cdot]$ is the Lie bracket.
 This is nonzero precisely because $\fg$ is nonabelian.
\item The degree-$4$ quartic obstruction class vanishes:
 $o_4(\ghat_k) = 0$. The vanishing is a consequence
 of the Jacobi identity for~$\fg$.
\item All higher obstruction classes vanish:
 $o_r(\ghat_k) = 0$ for $r \geq 4$.
\end{enumerate}
\end{theorem}

\begin{proof}
The degree-$2$ projection is computed in
Proposition~\ref{prop:kappa-km}.

For the cubic shadow, the degree-$3$ component of
$\Theta_{\ghat_k}$ lives in the cyclic deformation complex
$\mathrm{Def}_{\mathrm{cyc}}^{\mathrm{mod}}$ at degree~$3$.
The relevant graph sum reduces to a trivalent tree with
three external legs; the amplitude evaluates to the
structure constants $f^{abc}$ contracted with the
invariant form $(,)$, giving
$\mathfrak{C}(x,y,z) = (x,[y,z])$.
This vanishes if and only if $\fg$ is abelian.

The quartic obstruction $o_4$ lives in
$H^2(F^4\fg^{\mathrm{mod}}/F^5\fg^{\mathrm{mod}}, d_2)$.
The $d_2$ differential at degree~$4$ involves the graph
sum over trivalent trees with~$4$ external legs and one
internal edge. The obstruction is the class of
$[\Theta_3, \Theta_3]$ in the Lie bracket of the
convolution algebra.
The Jacobi identity
$f^{a[b}{}_e\, f^{c]d}{}_f + \text{cyclic} = 0$
forces this bracket to vanish in cohomology:
the quadri-linear form
$\sum_{e} f^{ab}{}_e\, f^{cd}{}_e$
satisfies the crossing symmetry
imposed by the Jacobi identity, and the resulting
cohomology class is trivially zero.
Since the quartic obstruction vanishes, all higher
obstructions vanish by induction: the master equation
$D\Theta + \tfrac{1}{2}[\Theta, \Theta] = 0$ at
degrees $r \geq 5$ reduces to $Do_r = 0$ with the
coboundary condition automatically satisfied.
\end{proof}

\begin{remark}[The collision residue]\label{rem:r-matrix}
The collision residue (binary genus-$0$ shadow)
of $V_k(\fg)$ is the Casimir $r$-matrix:
\begin{equation}\label{eq:km-rmatrix}
r(z) \;=\; \frac{\Omega}{(k + h^\vee)\, z},
\qquad
\Omega \;=\; \sum_a J^a \otimes J_a.
\end{equation}
The single pole reflects the fact that
the OPE~\eqref{eq:km-ope} has a double pole and the
$d\log$ kernel absorbs one power of $(z{-}w)$,
lowering the pole order by one.
This $r$-matrix satisfies the classical Yang--Baxter
equation.
\end{remark}

\begin{remark}[The critical discriminant]
\label{rem:discriminant}
The critical discriminant
$\Delta = 8\kappa S_4 = 0$ vanishes for all affine
Kac--Moody algebras (since $S_4 = 0$).
The shadow metric $Q_L(t) = (2\kappa + 3\alpha t)^2$
is a perfect square, consistent with finite shadow
depth. This is the algebraic characterization of
class~$L$: the shadow metric has a rational square root,
and the shadow connection has monodromy~$-1$ (the
Koszul sign).
\end{remark}

% ================================================================
% 7. EXAMPLES
% ================================================================

\section{Examples}\label{sec:examples}

\subsection{\texorpdfstring{$\widehat{\fsl}_2$}{sl2-hat}}

The simplest nonabelian case. Set $\fg = \fsl_2$,
$\dim\fg = 3$, $h^\vee = 2$, $\rk = 1$, $|\Delta_+| = 1$.

\begin{computation}\label{comp:sl2}
\begin{enumerate}[(i)]
\item \emph{Koszul dual level}:
 $k' = -k - 4$.
\item \emph{Modular characteristic}:
 $\kappa(\widehat{\fsl}_{2,k}) = 3(k+2)/4$.
\item \emph{Central charge}:
 $c(\widehat{\fsl}_{2,k}) = 3k/(k+2)$.
\item \emph{Complementarity}:
 $\kappa + \kappa' = 3(k{+}2)/4 + 3(-k{-}4{+}2)/4
 = 3(k{+}2)/4 + 3(-k{-}2)/4 = 0$. \checkmark
\item \emph{Bar cohomology}:
 $\dim H^n = 2n + 1$ for $n \geq 1$
 (Corollary~\ref{cor:bar-dim-formula}).
 Generating function $x(3{-}x)/(1{-}x)^2$.
\item \emph{$r$-matrix}:
 $r(z) = \Omega/((k{+}2)z)$, where
 $\Omega = e \otimes f + f \otimes e + h \otimes h/2$.
\item \emph{Shadow class}: $L$, $r_{\max} = 3$.
 Cubic shadow $\mathfrak{C}(x,y,z)
 = (x,[y,z])$ (the Killing form of $x$ with
 $[y,z]$); quartic vanishes by $\fsl_2$ Jacobi.
\end{enumerate}
\end{computation}

\begin{remark}\label{rem:sl2-admissible}
The simple quotient $L_k(\fsl_2)$ at admissible
levels $k = -2 + p/q$ (with $p \geq 2$, $q \geq 1$,
$\gcd(p,q) = 1$) is also chirally Koszul.
The argument uses the structural fact that
$L_k(\fsl_2)$ has a single family of null vectors
at each admissible level, and the bar cohomology is
controlled by the Kac--Wakimoto character formula.
\end{remark}

\subsection{\texorpdfstring{$\widehat{\fsl}_3$}{sl3-hat}}

Set $\fg = \fsl_3$, $\dim\fg = 8$, $h^\vee = 3$,
$\rk = 2$, $|\Delta_+| = 3$. The rank-$2$ case is
the first where a second independent Casimir appears
(the cubic Casimir $C_3 = d_{abc}\, J^a J^b J^c$),
and the Koszul dual structure acquires genuinely new
features relative to $\fsl_2$.

\begin{computation}\label{comp:sl3}
\begin{enumerate}[(i)]
\item \emph{Koszul dual level}: $k' = -k - 6$.
\item \emph{Modular characteristic}:
 $\kappa(\widehat{\fsl}_{3,k}) = 4(k{+}3)/3$.
 Verification: $\dim(\fsl_3) = 8$, $h^\vee = 3$, so
 $\kappa = 8(k{+}3)/(2 \cdot 3) = 4(k{+}3)/3$.
\item \emph{Central charge}:
 $c(\widehat{\fsl}_{3,k}) = 8k/(k{+}3)$.
\item \emph{Complementarity}:
 $\kappa(\widehat{\fsl}_{3,k}) + \kappa(\widehat{\fsl}_{3,k'})
 = \frac{4(k{+}3)}{3} + \frac{4(-k{-}6{+}3)}{3}
 = \frac{4(k{+}3) + 4(-k{-}3)}{3} = 0$.
\item \emph{Bar cohomology at low degree}:
 $\dim H^1 = 8$ (the adjoint representation).
 The total degree-$2$ bar cohomology has
 $\dim H^2 = 36$, computed by direct enumeration on
 the chiral bar complex
 $B^2 = \fg^{\otimes 2} \otimes \Omega^1_{\log}$
 (see~\cite{Lor-MKD}, Computation~7.2.1).
 This differs from the Chevalley--Eilenberg value
 $\dim H^2_{\CE}(\fn_-^{\mathrm{loop}}) = 20$
 because the chiral bar complex carries the
 Orlik--Solomon form factor, which is absent from
 the exterior algebra
 $\Lambda^*(\fn_-^{\mathrm{loop}})^*$.
 The generating function is conjectured to be
 \begin{equation}\label{eq:sl3-gf}
 P_{\widehat{\fsl}_3}(x)
 \;=\; 8x + 36 x^2 + \cdots.
 \end{equation}
\item \emph{Shadow class}: $L$, $r_{\max} = 3$.
 The cubic shadow $\mathfrak{C}_{\fsl_3}$ receives
 contributions from both the structure constants
 $f^{abc}$ and the symmetric tensor $d^{abc}$.
\item \emph{$r$-matrix}:
 $r(z) = \Omega_{\fsl_3}/((k{+}3)z)$, where
 $\Omega_{\fsl_3} = \sum_{a=1}^{8} T^a \otimes T_a$
 is the Casimir tensor in the Gell-Mann basis.
\item \emph{Koszul conductor}: $K = 2 \cdot 8 = 16$.
\end{enumerate}
\end{computation}

\begin{remark}\label{rem:sl3-vs-sl2}
The $\fsl_3$ case exhibits two qualitative differences
from $\fsl_2$.
First, the ring of $\fsl_3$-invariants in the symmetric
algebra on the adjoint is generated by Casimirs of
degrees~$2$ and~$3$ (rather than a single quadratic
Casimir for $\fsl_2$), producing faster-growing bar
cohomology.
Second, the cubic Casimir $d_{abc}$ provides a new
channel for the cubic shadow $\mathfrak{C}$ that does
not exist for $\fsl_2$ (where all cubic invariants
are proportional to the structure constants).
The shadow class remains $L$ because the quartic
obstruction still vanishes by the Jacobi identity,
which holds for all simple Lie algebras regardless
of rank.
\end{remark}

\subsection{Exceptional types}

\begin{proposition}\label{prop:exceptional}
For the simply-laced exceptional Lie algebras
$\fg \in \{E_6, E_7, E_8\}$, the invariants of
$V_k(\fg)$ are:
\begin{center}
\renewcommand{\arraystretch}{1.3}
\begin{tabular}{lccccc}
\toprule
$\fg$ & $\dim$ & $h^\vee$ & $\kappa(\ghat_k)$
 & $K$ & Class \\
\midrule
$E_6$ & $78$ & $12$ & $13(k{+}12)/4$ & $156$ & L \\
$E_7$ & $133$ & $18$ & $133(k{+}18)/36$ & $266$ & L \\
$E_8$ & $248$ & $30$ & $62(k{+}30)/15$ & $496$ & L \\
\bottomrule
\end{tabular}
\end{center}
Here $K = 2\dim\fg$ is the Koszul conductor.
Every entry satisfies $\kappa + \kappa' = 0$ and has
shadow depth $r_{\max} = 3$ \textup{(}class~$L$\textup{)}.
\end{proposition}

\begin{proof}
The formula $\kappa = \dim(\fg)(k{+}h^\vee)/(2h^\vee)$
is universal for all simple~$\fg$
(Proposition~\ref{prop:kappa-km}).
The Koszul conductor is $K = c + c'$; since
$c(V_k(\fg)) = k\,\dim\fg/(k{+}h^\vee)$ and
$c(V_{k'}(\fg)) = (-k{-}2h^\vee)\dim\fg/(-k{-}h^\vee)$,
we compute $K = \dim\fg[k/(k{+}h^\vee)
+ (-k{-}2h^\vee)/(-k{-}h^\vee)]$.
Simplifying:
\[
\frac{k}{k{+}h^\vee}
+ \frac{-k{-}2h^\vee}{-k{-}h^\vee}
= \frac{k}{k{+}h^\vee}
+ \frac{k{+}2h^\vee}{k{+}h^\vee}
= \frac{2(k{+}h^\vee)}{k{+}h^\vee} = 2.
\]
Therefore $K = 2\dim\fg$, universal for all simple~$\fg$.
Shadow class~$L$ follows from
Theorem~\ref{thm:km-class-L}.
\end{proof}

\begin{remark}[Non-simply-laced types]\label{rem:nsl}
For non-simply-laced $\fg \in \{B_n, C_n, G_2, F_4\}$,
the lacing number
$r^\vee = |\alpha_{\mathrm{long}}|^2/|\alpha_{\mathrm{short}}|^2$
enters the formulas only through the dual Coxeter
number~$h^\vee$ and dimension. The shadow depth remains
$r_{\max} = 3$ (class~$L$), the Koszul conductor is
still $K = 2\dim\fg$, and $\kappa + \kappa' = 0$.
All non-simply-laced affine Kac--Moody algebras are
chirally Koszul with the same structural properties as
the simply-laced cases. Specific values:
$\kappa(\widehat{G}_{2,k}) = 7(k{+}4)/4$,
$\kappa(\widehat{F}_{4,k}) = 26(k{+}9)/9$.
\end{remark}

% ================================================================
% 8. THE CRITICAL LEVEL
% ================================================================

\section{The critical level}\label{sec:critical}

The critical level $k = -h^\vee$ is the fixed point of the
Feigin--Frenkel involution $k \mapsto -k - 2h^\vee$ and a
singular point for both the Sugawara construction and the
modular characteristic. The bar complex simplifies
qualitatively.

\begin{proposition}\label{prop:critical-bar}
At the critical level $k = -h^\vee$:
\begin{enumerate}[(i)]
\item The curvature vanishes: $m_1^2 = 0$
 \textup{(}Proposition~\textup{\ref{prop:bar-curvature}}
 with $k + h^\vee = 0$\textup{)}.
 The bar complex is a strict dg coalgebra
 \textup{(}uncurved\textup{)}.
\item The modular characteristic vanishes:
 $\kappa(\ghat_{-h^\vee}) = 0$. The genus-$1$
 obstruction class $\mathrm{obs}_1 = 0$, and the
 partition function $F_1 = 0$.
\item The Sugawara construction is undefined: the
 denominator $k + h^\vee = 0$ in the Sugawara
 formula $L_{\mathrm{Sug}}
 = \frac{1}{2(k+h^\vee)} \sum_a J^a J^a$
 is singular. No Virasoro subalgebra exists.
\item The center of the completed enveloping algebra
 enlarges to the Feigin--Frenkel center
 $\mathfrak{z}(\ghat)$
 \textup{(}\cite{FF92}\textup{)},
 which is isomorphic to the algebra of functions on
 the space of $\fg^\vee$-opers on the formal disc.
\item The bar cohomology $H^*(\barB(V_{-h^\vee}(\fg)))$
 is related to the de Rham complex of the oper space:
 $H^n \cong \Omega^n(\mathrm{Op}_{\fg^\vee}(D))$
 for all $n \geq 0$
 \textup{(}via the fundamental local equivalence of
 Raskin~\cite{Raskin20}\textup{)}.
\end{enumerate}
\end{proposition}

\begin{proof}
Items~(i)--(ii) follow directly from
Proposition~\ref{prop:bar-curvature} and
equation~\eqref{eq:kappa-km} at $k = -h^\vee$.
Item~(iii) is classical~\cite{FF92}.
Item~(iv) is the Feigin--Frenkel theorem~\cite{FF92}.
Item~(v) follows from the bar-Ext identification
$H^n(\barB(V_k(\fg))) \cong
\Ext^n_{V_k(\fg)}(\mathbb{C}, \mathbb{C})$
combined with the fundamental local equivalence at
the critical level, which identifies
$\widehat{\fg}_{-h^\vee}$-modules with the Whittaker
category on $\mathrm{Op}_{\fg^\vee}(D)$.
\end{proof}

\begin{remark}[Critical level as self-dual point]
\label{rem:critical-self-dual}
The critical level is the unique self-dual point of
the Feigin--Frenkel involution:
$(-h^\vee)' = -(-h^\vee) - 2h^\vee = -h^\vee$.
The Koszul dual of the critical-level algebra is itself.
The vanishing of $\kappa$ at the critical level is
consistent with the self-duality:
$\kappa + \kappa' = 2\kappa = 0$.

In the shadow obstruction tower, $\kappa = 0$ means the
degree-$2$ projection vanishes. However, the cubic
shadow $\mathfrak{C}$ can still be nonzero (the Lie
bracket does not vanish at the critical level), so
the critical-level algebra is not class~$G$
(Gaussian). The classification is: at $k = -h^\vee$,
the algebra has $\kappa = 0$ with $\mathfrak{C} \neq 0$,
a degenerate class-$L$ configuration where the leading
scalar invariant vanishes but the cubic structure
persists.
\end{remark}

\begin{remark}[The geometric Langlands connection]
\label{rem:langlands}
The critical-level bar complex connects to the
geometric Langlands programme. The Feigin--Frenkel
center $\mathfrak{z}(\ghat) \cong
\mathrm{Fun}(\mathrm{Op}_{\fg^\vee}(D))$ is a
central ingredient of the geometric Langlands
proof~\cite{GLC24}.
The bar-Ext identification at the critical level
gives $H^n(\barB(V_{-h^\vee}(\fg)))
\cong \Omega^n(\mathrm{Op}_{\fg^\vee}(D))$,
so the bar cohomology of the critical-level vertex
algebra computes the algebraic de Rham cohomology
of the oper space. The local-to-global passage
(from the formal disc to a global curve~$X$)
transforms the bar complex into the factorization
algebra on $\mathrm{Ran}(X)$ whose global sections
recover the geometric Langlands kernel.
\end{remark}

% ================================================================
% 9. RELATION TO THE MONOGRAPH
% ================================================================

\section{Relation to the modular Koszul duality programme}
\label{sec:programme}

The results of this paper are proved as part of the
monograph~\cite{Lor-MKD}, where they appear in the
following form.

\begin{enumerate}[(1)]
\item Chiral Koszulness of $V_k(\fg)$
 (Theorem~\ref{thm:chiral-koszul-km}) is a special case
 of the PBW universality criterion
 (Proposition~5.3.1 of~\cite{Lor-MKD}):
 every freely strongly generated chiral algebra whose
 associated graded is classically Koszul is chirally Koszul.
 This criterion also applies to universal $\mathcal{W}$-algebras,
 lattice vertex algebras, and the $\beta\gamma$ system.

\item The Koszul dual identification
 (Theorem~\ref{thm:koszul-dual-km}) is a special case of
 the bar-cobar adjunction (Theorem~A of~\cite{Lor-MKD})
 combined with the Verdier intertwining theorem.

\item The modular characteristic formula
 (Proposition~\ref{prop:kappa-km}) and the complementarity
 sum $\kappa + \kappa' = 0$ are special cases of
 Theorem~D of~\cite{Lor-MKD}. The vanishing
 complementarity sum is specific to the Kac--Moody and
 free-field families; for $\mathcal{W}$-algebras,
 $\kappa + \kappa' = \varrho(\fg) \cdot K(\fg) \neq 0$,
 where $\varrho(\fg) = \sum_{i=1}^{r} (m_i + 1)^{-1}$ is
 the anomaly ratio
 (not the Weyl vector~$\rho$).

\item The shadow depth classification
 (Theorem~\ref{thm:km-class-L}) places affine Kac--Moody
 algebras in class~$L$ within the four-class partition
 G/L/C/M of~\cite{Lor-MKD}. The four classes are:
 Gaussian (Heisenberg, $r_{\max} = 2$),
 Lie/tree (affine KM, $r_{\max} = 3$),
 contact/quartic ($\beta\gamma$, $r_{\max} = 4$),
 and mixed (Virasoro, $\mathcal{W}_N$, $r_{\max} = \infty$).
 All four classes are chirally Koszul; shadow depth
 classifies complexity within the Koszul world,
 not Koszulness itself.

\item The collision residue $r(z) = \Omega/((k{+}h^\vee)z)$
 (Remark~\ref{rem:r-matrix}) is the binary genus-$0$
 projection of the universal Maurer--Cartan element
 $\Theta_{\ghat_k}$. Under the averaging map
 $\mathrm{av}\colon \fg^{E_1} \to \fg^{\mathrm{mod}}$,
 the affine modular characteristic is recovered with the
 Sugawara shift:
 $\mathrm{av}(r(z)) + \dim(\fg)/2 = \kappa(\ghat_k)$.
\end{enumerate}

% ================================================================
% BIBLIOGRAPHY
% ================================================================

\begin{thebibliography}{99}

\bibitem{BD04}
A.~Beilinson and V.~Drinfeld,
\emph{Chiral Algebras},
AMS Colloquium Publications, vol.~51,
American Mathematical Society, Providence, RI, 2004.

\bibitem{BGS96}
A.~Beilinson, V.~Ginzburg, and W.~Soergel,
Koszul duality patterns in representation theory,
\emph{J.~Amer.\ Math.\ Soc.}~\textbf{9} (1996), no.~2, 473--527.

\bibitem{FF92}
B.~Feigin and E.~Frenkel,
Affine Kac--Moody algebras at the critical level and
Gel'fand--Diki\u{\i} algebras,
\emph{Internat.\ J.\ Modern Phys.\ A}~\textbf{7}
(1992), suppl.~1A, 197--215.

\bibitem{FFR94}
B.~Feigin, E.~Frenkel, and N.~Reshetikhin,
Gaudin model, Bethe ansatz and critical level,
\emph{Comm.\ Math.\ Phys.}~\textbf{166} (1994), no.~1, 27--62.

\bibitem{FM94}
W.~Fulton and R.~MacPherson,
A compactification of configuration spaces,
\emph{Ann.\ of Math.~(2)}~\textbf{139} (1994), no.~1, 183--225.

\bibitem{GL76}
H.~Garland and J.~Lepowsky,
Lie algebra homology and the Macdonald--Kac formulas,
\emph{Invent.\ Math.}~\textbf{34} (1976), no.~1, 37--76.

\bibitem{Kostant61}
B.~Kostant,
Lie algebra cohomology and the generalized Borel--Weil theorem,
\emph{Ann.\ of Math.~(2)}~\textbf{74} (1961), no.~2, 329--387.

\bibitem{GLC24}
D.~Gaitsgory et al.,
Proof of the geometric Langlands conjecture,
2024, arXiv:2405.03599.

\bibitem{Lor-MKD}
R.~Lorgat,
\emph{Modular Koszul Duality for Chiral Algebras on Algebraic
Curves},
monograph in preparation, 2026.

\bibitem{Raskin20}
S.~Raskin,
Homological methods in semi-infinite contexts,
2020, arXiv:2002.01395.

\bibitem{LV12}
J.-L.~Loday and B.~Vallette,
\emph{Algebraic Operads},
Grundlehren der mathematischen Wissenschaften, vol.~346,
Springer, Berlin, 2012.

\bibitem{Priddy70}
S.~Priddy,
Koszul resolutions,
\emph{Trans.\ Amer.\ Math.\ Soc.}~\textbf{152} (1970),
39--60.

\end{thebibliography}

\end{document}
